% !TeX root = ../main.tex

\chapter{国内外研究现状}
\label{ch:related_work}

\section{通用目标检测方法演进}

通用目标检测方法的发展大致经历了 two-stage、one-stage、anchor-free 和 transformer 四类代表路线。two-stage 方法以 Faster R-CNN 为代表，通过候选区域生成与分类回归两阶段流程实现较高定位精度，在学术研究中长期占据重要地位\cite{ren2015_fasterrcnn}。其优点在于检测精度较高、定位能力较强，但推理速度和部署复杂度相对较高。

one-stage 方法以 YOLO 系列和 RetinaNet 为代表，直接进行密集预测，避免了显式候选框生成过程，因此在推理速度和工程部署方面具有明显优势\cite{redmon2016_yolo,redmon2018_yolov3,lin2017_focalloss}。这一路线更适合实时场景和资源受限环境，但在复杂背景与弱目标条件下容易出现漏检或误检问题。

anchor-free 方法以 FCOS 为代表，通过位置回归替代锚框设计，降低了 anchor 尺寸、比例和匹配策略带来的先验依赖\cite{tian2019_fcos}。该类方法简化了超参数设计，并在一定程度上提升了跨场景迁移能力，但其在微小目标场景下的优势仍需结合具体任务进一步验证。

transformer 检测器则通过全局关系建模增强上下文表达能力，典型方法包括 DETR、Deformable DETR 与 DINO\cite{carion2020_detr,zhu2021_deformabledetr,zhang2023_dino}。这类方法在端到端建模和全局建模方面具有较强潜力，但训练成本、数据需求和部署复杂度相对较高，不一定适合本科毕业设计周期中的完整复现实验。

总体来看，通用检测方法已经形成较成熟的技术谱系，不同路线在精度、速度、调参复杂度和部署难度之间具有明显差异。本文选择 DRENet、FCOS 和 YOLO26 进行对比，正是希望在专项方法、anchor-free 方法与 one-stage 方法之间建立较为清晰的横向参照。

\section{小目标检测研究进展}

小目标检测是目标检测研究中的一个重要分支，其难点主要体现在目标尺寸小、有效特征少、上下文依赖强和前景背景不平衡等方面。围绕这些问题，相关研究大致可以分为三个方向。

第一类方法关注\textbf{特征表达增强}。典型做法是利用 FPN 或其改进结构进行多尺度特征融合，使浅层细节与深层语义能够共同服务于小目标识别\cite{lin2017_fpn}。这类方法的优点是对现有检测框架兼容性较好，易于集成，但在纹理极弱的小目标场景下仍可能存在表达不足。

第二类方法关注\textbf{训练策略优化}。例如，Focal Loss 通过降低易分类样本的权重、强化难样本学习，缓解了密集检测中的类别不平衡问题\cite{lin2017_focalloss}。此外，数据增强、困难样本重采样和损失重加权等策略也常用于改善小目标训练效果。这一路线通常能够在不大幅修改模型结构的前提下带来收益，但其效果对数据质量和任务特性较为敏感。

第三类方法关注\textbf{后处理与推理策略改进}。在实际任务中，阈值设置、非极大值抑制以及多模型融合会显著影响 Precision 与 Recall 的平衡。对于目标稀疏、背景复杂的遥感场景而言，后处理策略往往不只是工程细节，而是影响最终表现的重要组成部分。

现有小目标检测研究为遥感微小船舶任务提供了重要启示，即需要同时考虑多尺度特征、样本不平衡和后处理策略。但由于遥感图像与自然场景在背景结构、目标分布和噪声形态上存在明显差异，通用小目标方法并不能直接保证在遥感场景中取得稳定效果，这也构成了本文进一步比较和分析的动机。

\section{遥感船舶检测研究现状}

遥感船舶检测是遥感目标检测的重要应用方向。随着公开数据集和开源工具链不断成熟，相关研究逐步从传统手工特征方法转向基于深度学习的端到端检测方法。LEVIR-Ship 数据集为中分辨率遥感微小船舶检测提供了较有代表性的公开基准，使研究者能够在统一数据条件下开展实验比较\cite{levirship2021_dataset}。

在具体方法上，DRENet 针对微小船舶特征弱和背景干扰强的问题，引入退化重建增强思想，在训练阶段强化目标表征能力，在 LEVIR-Ship 任务上表现出较好的检测效果\cite{chen2022_tinyship_drenet}。这一类专项方法说明，针对遥感微小目标进行结构或训练机制上的专门设计，能够在一定程度上弥补通用检测器的不足。

在工程生态方面，OpenMMLab 与 Ultralytics 等开源工具链为跨模型训练、测试和部署提供了成熟基础\cite{mmdet2018_openmmlab,ultralytics2024_docs}。这使得研究重点不再局限于单一模型精度，而更加关注对比协议是否统一、实验记录是否完整以及模型是否能够稳定接入实际系统。对于本科毕业设计而言，这种“算法实验与系统实现并重”的研究方式更具现实意义。

\begin{table}[htbp]
  \centering
  \caption{本文相关研究方向与代表路线}
  \label{tab:ch2_research_map}
  \begin{tabular}{lll}
    \toprule
    研究方向 & 代表路线 & 本文定位 \\
    \midrule
    通用检测 & Faster R-CNN / FCOS / YOLO / DETR & 作为跨范式对比基础 \\
    小目标检测 & 多尺度融合、损失重加权、增强策略 & 用于分析性能差异来源 \\
    遥感船舶检测 & LEVIR-Ship + 专项方法（如 DRENet） & 聚焦微小船舶场景 \\
    工程化实现 & mmdetection / ultralytics / 服务化封装 & 支撑系统落地与演示 \\
    \bottomrule
  \end{tabular}
\end{table}

\section{现有工作的不足与本文切入点}

综合上述研究可以看到，已有工作虽然为遥感微小船舶检测提供了丰富的方法基础，但仍存在以下不足。

首先，\textbf{实验协议不统一}。不同研究在数据划分、输入尺度、阈值设置和评测脚本上的差异，会导致结果难以进行严格横向比较。其次，\textbf{复现与工程细节披露不足}。部分方法仅给出模型结构和最终指标，却缺少足够完整的训练配置、日志留痕和部署说明，不利于毕业设计场景中的稳定复现。再次，\textbf{定性分析不充分}。对于复杂海况、近岸场景和纹理干扰导致的漏检、误检现象，已有工作往往停留在结果汇报层面，对差异来源分析不足。最后，\textbf{系统化闭环较弱}。从算法结果到可演示系统的完整链路，在不少研究中并未作为重点呈现。

针对这些不足，本文的切入点是：在统一协议下完成 DRENet、FCOS 和 YOLO26 的对比实验，用定量结果与定性样例分析不同范式的特点，并进一步构建轻量级检测系统，将算法结果转化为可交付、可展示和可回放的系统能力。这样的切入方式更符合本科毕业设计“研究论证与工程实现并重”的要求。

\section{本章小结}

本章系统梳理了通用目标检测、小目标检测和遥感船舶检测的研究进展，并总结了现有工作的主要不足。可以看到，当前研究不仅需要关注模型指标本身，还需要重视统一评测、过程复现和系统落地。基于这些认识，下一章将进一步介绍本文的数据集、评价指标、模型方案与统一实验协议。
